\subsection{Category of Iwahori--Whittaker perverse sheaves}
\label{jep104local034}

We now denote by 

$I^+ \subset G_{\jepPubScript{O}}$ the Iwahori subgroup associated with $B^+$. We also denote by $I^+_\mathrm{u}$ the pro-unipotent radical of $I^+$, i.e.~the inverse image of $U^+$ under the map $I^+ \to B^+$.

We assume that there exists a primitive $p$-th root of unity in $\Bbbk$, and fix one. This choice determines a
character $\psi$ of the prime subfield of $\jepPubBbb{F}$ (with values in $\Bbbk^\times$), and we denote by $\jepPubEuler{L}^\Bbbk_\psi$ the corresponding Artin--Schreier local system on $\jepPubBbb{G}_{\mathrm{a}}$. (Below, some arguments using Verdier duality will also involve the Artin--Schreier local system $\jepPubEuler{L}^\Bbbk_{-\psi}$ associated with the character $\psi^{-1}$; clearly these two versions play similar roles.) We also consider the ``generic'' character $\chi : U^+ \to \jepPubBbb{G}_{\mathrm{a}}$ 

defined as the composition
\[
 U^+ \twoheadrightarrow U^+ / [U^+,U^+] \xleftarrow[\sim]{\prod_{\alpha} u_\alpha} \prod_{\alpha \in \Delta_{\mathrm{s}}} \jepPubBbb{G}_{\mathrm{a}} \xrightarrow{+} \jepPubBbb{G}_{\mathrm{a}},
\]
and denote by $\chi_{I^+}$ its composition with the projection $I^+_\mathrm{u} \twoheadrightarrow U^+$. We can then define the ``Iwahori--Whittaker'' derived category
\[
D^{\mathrm{b}}_{\jepPubEuler{IW}}(\mathrm{Gr}, \Bbbk)
\]

as the $(I^+_\mathrm{u}, \chi_{I^+}^*(\jepPubEuler{L}^\Bbbk_\psi))$-equivariant constructible derived category of $\Bbbk$-sheaves on $\mathrm{Gr}$ (see e.g.~\cite[App.~A]{jep104ref03} for a review of the construction of this category). This category admits a perverse t-structure, whose heart will be denoted $\mathrm{Perv}_{\jepPubEuler{IW}}(\mathrm{Gr},\Bbbk)$, and moreover the ``realization functor''
\[
D^{\mathrm{b}} \mathrm{Perv}_{\jepPubEuler{IW}}(\mathrm{Gr},\Bbbk) \to D^{\mathrm{b}}_{\jepPubEuler{IW}}(\mathrm{Gr}, \Bbbk)
\]
is an equivalence of triangulated categories.

Note that any ring of coefficients considered above appears in an $\ell$-modular triple $(\jepPubBbb{K},\jepPubBbb{O},\jepPubBbb{L})$ where $\jepPubBbb{K}$ is a finite extension of $\jepPubBbb{Q}_\ell$, $\jepPubBbb{O}$ is its ring of integers, and $\jepPubBbb{L}$ is the residue field of $\jepPubBbb{O}$. In this setting the embedding $\jepPubBbb{O}\hookrightarrow\jepPubBbb{K}$ and the projection $\jepPubBbb{O}\to\jepPubBbb{L}$ induce bijections between the $p$-th roots of unity in $\jepPubBbb{O}$, $\jepPubBbb{K}$ and $\jepPubBbb{L}$. Therefore, choosing a primitive root in any of these rings provides primitive roots in all three rings, and we can then consider extension of scalars functors
\[
\jepPubBbb{K}\jepDerivedTensor_{\jepPubBbb{O}}(-):D^{\mathrm{b}}_{\jepPubEuler{IW}}(\mathrm{Gr},\jepPubBbb{O})\to D^{\mathrm{b}}_{\jepPubEuler{IW}}(\mathrm{Gr},\jepPubBbb{K}),\qquad
\jepPubBbb{L}\jepDerivedTensor_{\jepPubBbb{O}}(-):D^{\mathrm{b}}_{\jepPubEuler{IW}}(\mathrm{Gr},\jepPubBbb{O})\to D^{\mathrm{b}}_{\jepPubEuler{IW}}(\mathrm{Gr},\jepPubBbb{L}).
\]
These functors will play a crucial role in our arguments below. (Here the first functor is $t$-exact, but the second one is only right $t$-exact.)
For $\lambda \in \boldsymbol{X}^\vee$ we set
\[
X_\lambda := I^+ \cdot L_\lambda.
\]
Then again we have
\[
\mathrm{Gr} = \mathop{\textstyle\bigsqcup}\limits_{\lambda \in \boldsymbol{X}^\vee} X_\lambda.
\]

\begin{lem}
\label{jep104local020}
The orbit $X_\lambda$ supports an $(I^+_\mathrm{u}, \chi_{I^+}^*(\jepPubEuler{L}^\Bbbk_\psi))$-equivariant local system iff $\lambda \in \boldsymbol{X}^\vee_{++}$.
\end{lem}

\begin{proof}[Sketch of proof]

Let $\lambda \in \boldsymbol{X}^\vee_+$, and consider the affine bundle $\mathrm{Gr}^\lambda \to G/P_\lambda$ (see~Section~\ref{jep104local030}).
The decomposition of $\mathrm{Gr}^\lambda$ in $I_\mathrm{u}^+$-orbits is obtained by pullback from the decomposition of $G/P_\lambda$ into $U^+$-orbits; in particular, for $\mu\in W_{\mathrm{f}}\cdot\lambda$, $X_\mu$ supports an $(I^+_\mathrm{u}, \chi_{I^+}^*(\jepPubEuler{L}^\Bbbk_\psi))$-equivariant local system iff its image in $G/P_\lambda$ is a free $U^+$-orbit. If $\lambda \notin \boldsymbol{X}^\vee_{++}$ there is no such orbit in $G/P_\lambda$, and if $\lambda \in \boldsymbol{X}^\vee_{++}$ there is exactly one, corresponding to $X_\lambda$.
\end{proof}

Note that if $\lambda \in \boldsymbol{X}^\vee_{++}$, since $X_\lambda$ is open dense in $\mathrm{Gr}^\lambda$ we have
\begin{equation}
\label{jep104local007}
\dim(X_\lambda) = \langle \lambda, 2\rho \rangle
\end{equation}
and
\[
X_\lambda \subset \overline{X_\mu} \quad \text{iff} \quad \lambda \preccurlyeq \mu.
\]
For $\lambda \in \boldsymbol{X}^\vee_{++}$ we will denote by
\[
\Delta^{\jepPubEuler{IW}}_\lambda(\Bbbk), \quad \text{resp.} \quad \nabla^{\jepPubEuler{IW}}_\lambda(\Bbbk),
\]
the standard, resp.~costandard, $(I^+_\mathrm{u}, \chi_{I^+}^*(\jepPubEuler{L}^\Bbbk_\psi))$-equivariant perverse sheaf on $\mathrm{Gr}$ associated with $\lambda$, i.e.~the $!$-extension, resp.~$*$-extension, to $\mathrm{Gr}$ of the free rank-$1$ $(I^+_\mathrm{u}, \chi_{I^+}^*(\jepPubEuler{L}^\Bbbk_\psi))$-equivariant perverse sheaf on $X_\lambda$. (Once again, these objects are perverse sheaves thanks to~\cite[Cor.~4.1.3]{jep104ref12}.) We will also denote by $\mathrm{IC}_\lambda^{\jepPubEuler{IW}}(\Bbbk)$ the image of any generator of the rank-$1$ free $\Bbbk$-module
\[
\operatorname{Hom}_{\mathrm{Perv}_{\jepPubEuler{IW}}(\mathrm{Gr}, \Bbbk)}(\Delta_\lambda^{\jepPubEuler{IW}}(\Bbbk), \nabla_\lambda^{\jepPubEuler{IW}}(\Bbbk)).
\]
If $\Bbbk$ is a field then $\mathrm{IC}_\lambda^{\jepPubEuler{IW}}(\Bbbk)$ is a simple perverse sheaf.

Note that since $\varsigma$ is minimal in $\boldsymbol{X}^\vee_{++}$ for $\preccurlyeq$, we have
\begin{equation}
\label{jep104local003}
\Delta^{\jepPubEuler{IW}}_{\varsigma}(\Bbbk) = \nabla^{\jepPubEuler{IW}}_\varsigma(\Bbbk) = \mathrm{IC}^{\jepPubEuler{IW}}_\varsigma(\Bbbk).
\end{equation}

\begin{lem}
\label{jep104local021}
Assume that $\Bbbk$ is a field of characteristic $0$. Then the $i$-th cohomology of the stalks of $\mathrm{IC}^{\jepPubEuler{IW}}_\lambda(\Bbbk)$ vanish unless $i \equiv \dim X_\lambda \pmod 2$.
\end{lem}

\begin{proof}[Sketch of proof]
Since the morphism $\pi$ is smooth, by standard properties of perverse sheaves (see e.g.~\cite[\S 4.2.6]{jep104ref12}) it suffices to prove a similar statement on $\mathrm{Fl}$ instead of $\mathrm{Gr}$. Now the Decomposition Theorem implies that all simple $(I_\mathrm{u}^+, \chi_{I^+}^*(\jepPubEuler{L}^\Bbbk_\psi))$-equivariant perverse sheaves on $\mathrm{Fl}$ can be obtained from the one corresponding to the orbit of the base point by convolving on the right with $I^-$-equivariant simple perverse sheaves on $\mathrm{Fl}$ corresponding to orbits of dimension either $0$ or $1$. Standard arguments (going back at least to~\cite{jep104ref43}) show that these operations preserve the parity-vanishing property of stalks, and the claim follows.
\end{proof}

\begin{rmk}
\label{jep104local027}
Assume that $\Bbbk$ is a field.
Following~\cite{jep104ref31}\footnote{In~\cite{jep104ref31} the authors consider the setting of ``ordinary'' constructible complexes. However, as observed already in~\cite[\S 11.1]{jep104ref42} or~\cite[\S 6.2]{jep104ref02}, their considerations apply verbatim in our Iwahori--Whittaker setting.} we will say that an object of $D^{\mathrm{b}}_{\jepPubEuler{IW}}(\mathrm{Gr}, \Bbbk)$ is \emph{even}, resp.~\emph{odd}, if its restriction and corestriction to each stratum is concentrated in even, resp.~odd, degrees, and that it is \emph{parity} if it is isomorphic to a direct sum $\jepPubEuler{F} \oplus \jepPubEuler{F}'$ with $\jepPubEuler{F}$ even and $\jepPubEuler{F}'$ odd. Using this language, Lemma~\ref{jep104local021} states that if $\mathrm{char}(\Bbbk)=0$ then the objects $\mathrm{IC}^{\jepPubEuler{IW}}_\lambda(\Bbbk)$ are parity, of the same parity as $\dim(X_\lambda)$.
\end{rmk}

\begin{cor}
\label{jep104local001}
Assume that $\Bbbk$ is a field.
The category $\mathrm{Perv}_{\jepPubEuler{IW}}(\mathrm{Gr}, \Bbbk)$ is a highest weight category with weight poset $(\boldsymbol{X}^\vee_{++}, \preccurlyeq)$, standard objects $\{\Delta_\lambda^{\jepPubEuler{IW}}(\Bbbk) : \lambda \in \boldsymbol{X}^\vee_{++}\}$, and costandard objects $\{\nabla^{\jepPubEuler{IW}}_\lambda(\Bbbk) : \lambda \in \boldsymbol{X}^\vee_{++}\}$. Moreover, if $\mathrm{char}(\Bbbk)=0$ then this category is semisimple.
\end{cor}

\begin{proof}
The first claim is standard, as e.g.~in~\cite[\S 3.3]{jep104ref14}. For the second claim, we observe that the orbits $X_\lambda$ (for $\lambda\in\boldsymbol{X}^\vee_{++}$) have dimensions of constant parity on each connected component of $\mathrm{Gr}$, see~\eqref{jep104local007}. 

Using this and Lemma~\ref{jep104local021},
the semisimplicity can be proved exactly as in the case of the category $\mathrm{Perv}_{G_{\jepPubScript{O}}}(\mathrm{Gr}, \Bbbk)$. Namely, we have to prove that
\[
\operatorname{Ext}^1_{\mathrm{Perv}_{\jepPubEuler{IW}}(\mathrm{Gr}, \Bbbk)}(\mathrm{IC}^{\jepPubEuler{IW}}_\lambda(\Bbbk),\mathrm{IC}^{\jepPubEuler{IW}}_\mu(\Bbbk))=\operatorname{Hom}_{D^{\mathrm{b}}_{\jepPubEuler{IW}}(\mathrm{Gr}, \Bbbk)}(\mathrm{IC}^{\jepPubEuler{IW}}_\lambda(\Bbbk),\mathrm{IC}^{\jepPubEuler{IW}}_\mu(\Bbbk)[1])
\]
vanishes for any $\lambda,\mu$. If $X_\lambda$ and $X_\mu$ belong to different connected components of $\mathrm{Gr}$ then this claim is obvious; otherwise $\mathrm{IC}^{\jepPubEuler{IW}}_\lambda(\Bbbk)$ and $\mathrm{IC}^{\jepPubEuler{IW}}_\mu(\Bbbk)$ are either both even or both odd (see Remark~\ref{jep104local027}), so that the desired vanishing follows from~\cite[Cor.~2.8]{jep104ref31}.

\end{proof}

\begin{rmk}\phantomsection
\label{jep104local025}
\begin{enumerate}
\item
\label{jep104local013}
Once Corollary~\ref{jep104local001} is known, one can refine Lemma~\ref{jep104local021} drastically: if $\Bbbk$ is a field of characteristic $0$, then the simple perverse sheaves $\mathrm{IC}^{\jepPubEuler{IW}}_\lambda(\Bbbk)$ are clean, in the sense that if $i_\mu : X_\mu \to \mathrm{Gr}$ is the embedding, for any $\mu \neq \lambda$ we have
\begin{equation}
\label{jep104local004}
i_\mu^* \bigl( \mathrm{IC}^{\jepPubEuler{IW}}_\lambda(\Bbbk) \bigr) = i_\mu^! \bigl( \mathrm{IC}^{\jepPubEuler{IW}}_\lambda(\Bbbk) \bigr)=0.
\end{equation}
In fact, as in Remark~\ref{jep104local024}\eqref{jep104local014}, the semisimplicity claim in Corollary~\ref{jep104local001} implies that the natural maps
\[
 \Delta^{\jepPubEuler{IW}}_\lambda(\Bbbk) \to \mathrm{IC}^{\jepPubEuler{IW}}_\lambda(\Bbbk) \to \nabla^{\jepPubEuler{IW}}_\lambda(\Bbbk)
\]
are isomorphisms, which is equivalent to~\eqref{jep104local004}.
(See also~\cite[Cor.~2.2.3]{jep104ref09} for a different proof of~\eqref{jep104local004}.)
This observation can be used to give a new proof of the main result of~\cite{jep104ref26}, hence of the geometric Casselman--Shalika formula.

\item
\label{jep104local012}
The same arguments as in Remark~\ref{jep104local024}\eqref{jep104local011} show that for any coefficients $\Bbbk$, any $\lambda,\mu \in \boldsymbol{X}^\vee_{++}$ and any $n \in \jepPubBbb{Z}$ we have
 \[
  \operatorname{Hom}_{D^{\mathrm{b}} \mathrm{Perv}_{\jepPubEuler{IW}}(\mathrm{Gr}, \Bbbk)} \left( \Delta_\lambda^{\jepPubEuler{IW}}(\Bbbk), \nabla_\mu^{\jepPubEuler{IW}}(\Bbbk)[n] \right) =
  \begin{cases}
   \Bbbk & \text{if $\lambda=\mu$ and $n=0$;} \\
   0 & \text{otherwise.}
  \end{cases}
 \]
 (In this case, the existence of enough projectives in $\mathrm{Perv}_{\jepPubEuler{IW}}(Z,\Bbbk)$ can be checked using the techniques of~\cite[\S 2]{jep104ref41}.)
\end{enumerate}
\end{rmk}

